\section{Continuous Assessment 1 (CA-1)}

\ExamHeader{Continuous Assessment 1 (CA-1)}{Set B (Model Paper 2)}{60 Minutes}{30 Marks}{None}

\noindent
\textbf{General Instructions:}
\begin{enumerate}[topsep=2pt,itemsep=1pt]
    \item This paper contains \textbf{6 subjective questions} carrying \textbf{5 marks each}.
    \item Candidates are required to \textbf{attempt any 5 questions} ($5 \times 5 = 25$ marks or scaled to 30 marks).
    \item Show all necessary algebraic steps, matrix reductions, and derivative formulations clearly.
\end{enumerate}
\vspace{3mm}

\noindent\textbf{Question 1} \hfill \textbf{[5 Marks]}

Determine the rank of the matrix $A = \begin{bmatrix} 1 & 2 & -1 & 3 \\ 2 & 4 & 1 & -2 \\ 3 & 6 & 3 & -7 \end{bmatrix}$ by reducing it to Row Echelon Form.

\begin{solbox}
Apply elementary row operations:
\begin{enumerate}
    \item $R_2 \to R_2 - 2R_1 \implies \begin{bmatrix} 1 & 2 & -1 & 3 \\ 0 & 0 & 3 & -8 \\ 3 & 6 & 3 & -7 \end{bmatrix}$
    \item $R_3 \to R_3 - 3R_1 \implies \begin{bmatrix} 1 & 2 & -1 & 3 \\ 0 & 0 & 3 & -8 \\ 0 & 0 & 6 & -16 \end{bmatrix}$
    \item $R_3 \to R_3 - 2R_2 \implies \begin{bmatrix} 1 & 2 & -1 & 3 \\ 0 & 0 & 3 & -8 \\ 0 & 0 & 0 & 0 \end{bmatrix}$
\end{enumerate}
The row echelon form contains \textbf{2 non-zero rows}. Hence, $\operatorname{rank}(A) = 2$.
\end{solbox}

\begin{examinertip}
State each row operation explicitly. The rank equals the number of non-zero rows in the final echelon form.
\end{examinertip}

\vspace{4mm}

\noindent\textbf{Question 2} \hfill \textbf{[5 Marks]}

Find the values of $\lambda$ and $\mu$ such that the system of equations $x+y+z=6$, $x+2y+3z=10$, $x+2y+\lambda z=\mu$ has (i) no solution, (ii) a unique solution, and (iii) infinitely many solutions.

\begin{solbox}
Write the augmented matrix $[A|b]$:
\begin{equation}
\begin{bmatrix} 1 & 1 & 1 & 6 \\ 1 & 2 & 3 & 10 \\ 1 & 2 & \lambda & \mu \end{bmatrix} \sim \begin{bmatrix} 1 & 1 & 1 & 6 \\ 0 & 1 & 2 & 4 \\ 0 & 0 & \lambda-3 & \mu-10 \end{bmatrix}
\end{equation}
\begin{itemize}
    \item \textbf{No Solution}: $\lambda = 3$ and $\mu \ne 10$ ($\operatorname{rank}(A)=2 \ne \operatorname{rank}([A|b])=3$).
    \item \textbf{Unique Solution}: $\lambda \ne 3$ for any real $\mu$ ($\operatorname{rank}(A)=\operatorname{rank}([A|b])=3$).
    \item \textbf{Infinite Solutions}: $\lambda = 3$ and $\mu = 10$ ($\operatorname{rank}(A)=\operatorname{rank}([A|b])=2 < 3$).
\end{itemize}
\end{solbox}

\begin{examinertip}
Apply Rouch\'e--Capelli theorem directly to the pivot entries $(\lambda-3)$ and $(\mu-10)$.
\end{examinertip}

\vspace{4mm}

\noindent\textbf{Question 3} \hfill \textbf{[5 Marks]}

Find the eigenvalues and eigenvectors of $A = \begin{bmatrix} 3 & 2 \\ 2 & 3 \end{bmatrix}$.

\begin{solbox}
Characteristic equation: $|A - \lambda I| = (3-\lambda)^2 - 4 = \lambda^2 - 6\lambda + 5 = 0 \implies (\lambda-1)(\lambda-5)=0$.
Eigenvalues: $\lambda_1 = 5, \; \lambda_2 = 1$.
\begin{itemize}
    \item For $\lambda = 5$: $\begin{bmatrix} -2 & 2 \\ 2 & -2 \end{bmatrix}\begin{bmatrix} x \\ y \end{bmatrix} = 0 \implies x = y \implies v_1 = \begin{bmatrix} 1 \\ 1 \end{bmatrix}$.
    \item For $\lambda = 1$: $\begin{bmatrix} 2 & 2 \\ 2 & 2 \end{bmatrix}\begin{bmatrix} x \\ y \end{bmatrix} = 0 \implies x = -y \implies v_2 = \begin{bmatrix} 1 \\ -1 \end{bmatrix}$.
\end{itemize}
\end{solbox}

\begin{examinertip}
Notice $v_1 \cdot v_2 = 1(1) + 1(-1) = 0$, verifying orthogonality for symmetric matrices.
\end{examinertip}

\vspace{4mm}

\noindent\textbf{Question 4} \hfill \textbf{[5 Marks]}

If $y = e^{a \sin^{-1}x}$, prove that $(1-x^2)y_{n+2} - (2n+1)x y_{n+1} - (n^2+a^2)y_n = 0$.

\begin{solbox}
1. First derivative: $y_1 = \frac{a y}{\sqrt{1-x^2}} \implies (1-x^2)y_1^2 = a^2 y^2$.
2. Second derivative: $(1-x^2)(2y_1 y_2) - 2x y_1^2 = 2a^2 y y_1 \implies (1-x^2)y_2 - x y_1 - a^2 y = 0$.
3. Differentiate $n$ times using Leibniz theorem:
\begin{align}
[(1-x^2)y_{n+2} - 2nx y_{n+1} - n(n-1)y_n] - [x y_{n+1} + n y_n] - a^2 y_n &= 0 \\
(1-x^2)y_{n+2} - (2n+1)x y_{n+1} - (n^2+a^2)y_n &= 0
\end{align}
\end{solbox}

\begin{examinertip}
Differentiating the squared identity $(1-x^2)y_1^2 = a^2 y^2$ saves substantial time over quotient rule.
\end{examinertip}

\vspace{4mm}

\noindent\textbf{Question 5} \hfill \textbf{[5 Marks]}

Verify Rolle's Theorem for $f(x) = x^3 - 4x$ on the interval $[-2, 2]$.

\begin{solbox}
1. $f(x)$ is a polynomial, continuous on $[-2, 2]$ and differentiable on $(-2, 2)$.
2. $f(-2) = (-2)^3 - 4(-2) = 0$ and $f(2) = 2^3 - 4(2) = 0 \implies f(-2) = f(2)$.
3. Setting $f'(c) = 3c^2 - 4 = 0 \implies c = \pm \frac{2}{\sqrt{3}} \approx \pm 1.155$.
Since both values lie strictly in $(-2, 2)$, Rolle's theorem is verified.
\end{solbox}

\begin{examinertip}
Check all three hypotheses before computing $f'(c) = 0$.
\end{examinertip}

\vspace{4mm}

\noindent\textbf{Question 6} \hfill \textbf{[5 Marks]}

Evaluate $\lim_{x\to 0} \left(\frac{\sin x}{x}\right)^{1/x^2}$ using logarithmic transformation and L'H\^opital's Rule.

\begin{solbox}
This is of indeterminate form $[1^\infty]$.
Let $y = \left(\frac{\sin x}{x}\right)^{1/x^2} \implies \ln y = \frac{\ln(\sin x / x)}{x^2}$.
As $x \to 0$, this is of form $\left[\frac{0}{0}\right]$.
Applying L'H\^opital's rule:
\begin{equation}
\lim_{x\to 0} \frac{\frac{x}{\sin x} \cdot \frac{x\cos x - \sin x}{x^2}}{2x} = \lim_{x\to 0} \frac{x\cos x - \sin x}{2x^2 \sin x}
\end{equation}
Using series expansion $\sin x = x - x^3/6 + \dots, \cos x = 1 - x^2/2 + \dots$:
\begin{equation}
x(1 - x^2/2) - (x - x^3/6) = -x^3/3 \implies \lim_{x\to 0} \frac{-x^3/3}{2x^3} = -\frac{1}{6}
\end{equation}
Thus $\lim y = e^{-1/6} = \frac{1}{\sqrt[6]{e}}$.
\end{solbox}

\begin{examinertip}
Maclaurin series expansion of numerator eliminates repeated tedious applications of L'H\^opital's rule.
\end{examinertip}

\vspace{4mm}

